\textbf{Phần 1. Trắc nghiệm lựa chọn}
\setcounterpageref{ex}{0}

\begin{ex} %[3H3V53]
Một hình chữ nhật có chiều dài $16\text{ cm}$, chiều rộng bằng một nửa chiều dài. Diện tích của hình chữ nhật đó là
\choice
{$64\text{ cm}^2$}
{$96\text{ cm}^2$}
{\True $128\text{ cm}^2$}
{$48\text{ cm}^2$}
\loigiai{
Chiều rộng: $16 : 2 = 8\text{ cm}$. Diện tích: $16 \times 8 = 128\text{ (cm}^2)$.
\textbf{Đáp án của câu hỏi là C}
}
\end{ex}

\begin{ex} %[3H3V53]
Một hình vuông có diện tích $36\text{ cm}^2$. Chu vi của hình vuông đó là
\choice
{\True $24\text{ cm}$}
{$18\text{ cm}$}
{$36\text{ cm}$}
{$12\text{ cm}$}
\loigiai{
Vì $6 \times 6 = 36$ nên cạnh hình vuông là $6\text{ cm}$. Chu vi: $6 \times 4 = 24\text{ (cm)}$.
\textbf{Đáp án của câu hỏi là A}
}
\end{ex}

\begin{ex} %[3H3V53]
Một mảnh đất hình chữ nhật có chu vi $40\text{ m}$, chiều dài hơn chiều rộng $4\text{ m}$. Diện tích của mảnh đất đó là
\choice
{$80\text{ m}^2$}
{$100\text{ m}^2$}
{$84\text{ m}^2$}
{\True $96\text{ m}^2$}
\loigiai{
Nửa chu vi: $40 : 2 = 20\text{ m}$. Chiều dài: $(20 + 4) : 2 = 12\text{ m}$. Chiều rộng: $12 - 4 = 8\text{ m}$. Diện tích: $12 \times 8 = 96\text{ (m}^2)$.
\textbf{Đáp án của câu hỏi là D}
}
\end{ex}

\begin{ex} %[3H3V53]
Người ta rào một mảnh vườn hình chữ nhật có chiều dài $30\text{ m}$ và chiều rộng $18\text{ m}$. Cọc rào cách nhau $2\text{ m}$. Cần tất cả bao nhiêu cọc rào?
\choice
{$46$ cọc}
{\True $48$ cọc}
{$50$ cọc}
{$24$ cọc}
\loigiai{
Chu vi vườn: $(30 + 18) \times 2 = 96\text{ m}$. Số cọc rào khép kín: $96 : 2 = 48$ cọc.
\textbf{Đáp án của câu hỏi là B}
}
\end{ex}

\begin{ex} %[3H3V53]
Cắt một tấm bìa hình vuông cạnh $10\text{ cm}$ thành 4 hình vuông nhỏ bằng nhau. Diện tích của mỗi hình vuông nhỏ là
\choice
{$20\text{ cm}^2$}
{$50\text{ cm}^2$}
{\True $25\text{ cm}^2$}
{$16\text{ cm}^2$}
\loigiai{
Diện tích hình vuông lớn: $10 \times 10 = 100\text{ cm}^2$. Mỗi hình nhỏ: $100 : 4 = 25\text{ (cm}^2)$.
\textbf{Đáp án của câu hỏi là C}
}
\end{ex}

\begin{ex} %[3H3V53]
Một hình chữ nhật có chiều dài $15\text{ cm}$ và diện tích $90\text{ cm}^2$. Chu vi hình chữ nhật đó là
\choice
{\True $42\text{ cm}$}
{$30\text{ cm}$}
{$45\text{ cm}$}
{$36\text{ cm}$}
\loigiai{
Chiều rộng: $90 : 15 = 6\text{ cm}$. Chu vi: $(15 + 6) \times 2 = 42\text{ (cm)}$.
\textbf{Đáp án của câu hỏi là A}
}
\end{ex}

\begin{ex} %[3H3V53]
Một mảnh vườn hình chữ nhật có chiều dài $25\text{ m}$. Nếu bớt chiều dài đi $5\text{ m}$ thì mảnh vườn trở thành hình vuông. Diện tích của mảnh vườn ban đầu là
\choice
{$400\text{ m}^2$}
{$600\text{ m}^2$}
{$450\text{ m}^2$}
{\True $500\text{ m}^2$}
\loigiai{
Cạnh hình vuông bằng chiều rộng mảnh vườn: $25 - 5 = 20\text{ m}$. Diện tích ban đầu: $25 \times 20 = 500\text{ (m}^2)$.
\textbf{Đáp án của câu hỏi là D}
}
\end{ex}

\begin{ex} %[3H3V53]
Cho hình vuông có chu vi $48\text{ cm}$. Chia hình vuông đó thành 2 hình chữ nhật bằng nhau. Chu vi của mỗi hình chữ nhật là
\choice
{$24\text{ cm}$}
{\True $36\text{ cm}$}
{$30\text{ cm}$}
{$40\text{ cm}$}
\loigiai{
Cạnh hình vuông: $48 : 4 = 12\text{ cm}$. Mỗi hình chữ nhật có kích thước $12\text{ cm} \times 6\text{ cm}$. Chu vi: $(12 + 6) \times 2 = 36\text{ (cm)}$.
\textbf{Đáp án của câu hỏi là B}
}
\end{ex}

\begin{ex} %[3H3V53]
Một tam giác có chu vi $60\text{ cm}$. Cạnh thứ nhất dài $18\text{ cm}$, cạnh thứ hai dài $22\text{ cm}$. Cạnh thứ ba dài là
\choice
{$15\text{ cm}$}
{$25\text{ cm}$}
{\True $20\text{ cm}$}
{$24\text{ cm}$}
\loigiai{
Cạnh thứ ba: $60 - (18 + 22) = 20\text{ (cm)}$.
\textbf{Đáp án của câu hỏi là C}
}
\end{ex}

\begin{ex} %[3H3V53]
Một bức tường hình chữ nhật có chiều dài $6\text{ m}$, chiều cao $3\text{ m}$. Trên tường có một cửa sổ diện tích $2\text{ m}^2$. Diện tích cần sơn tường là
\choice
{\True $16\text{ m}^2$}
{$18\text{ m}^2$}
{$14\text{ m}^2$}
{$20\text{ m}^2$}
\loigiai{
Diện tích cả tường: $6 \times 3 = 18\text{ m}^2$. Diện tích cần sơn: $18 - 2 = 16\text{ (m}^2)$.
\textbf{Đáp án của câu hỏi là A}
}
\end{ex}

\begin{ex} %[3H3V53]
Hình vuông $A$ có cạnh gấp 3 lần cạnh hình vuông $B$. Chu vi hình vuông $A$ gấp mấy lần chu vi hình vuông $B$?
\choice
{$9$ lần}
{$6$ lần}
{$12$ lần}
{\True $3$ lần}
\loigiai{
Chu vi tỉ lệ thuận với cạnh: cạnh gấp 3 lần thì chu vi cũng gấp 3 lần.
\textbf{Đáp án của câu hỏi là D}
}
\end{ex}

\begin{ex} %[3H3V53]
Một hình chữ nhật có chu vi $32\text{ cm}$. Chiều dài gấp 3 lần chiều rộng. Diện tích của hình chữ nhật đó là
\choice
{$64\text{ cm}^2$}
{\True $48\text{ cm}^2$}
{$36\text{ cm}^2$}
{$56\text{ cm}^2$}
\loigiai{
Nửa chu vi: $32 : 2 = 16\text{ cm}$. Rộng: $16 : (3 + 1) = 4\text{ cm}$. Dài: $4 \times 3 = 12\text{ cm}$. Diện tích: $12 \times 4 = 48\text{ (cm}^2)$.
\textbf{Đáp án của câu hỏi là B}
}
\end{ex}

\textbf{Phần 2. Trắc nghiệm đúng sai}
\setcounterpageref{ex}{0}

\begin{ex} %[3H3V53]
Xét tính đúng sai của các nhận định ôn tập chu vi và diện tích:
\choiceTF
{\True Đơn vị đo chu vi là đơn vị độ dài ($\text{cm}, \text{m}$), đơn vị diện tích là đơn vị diện tích ($\text{cm}^2, \text{m}^2$).}
{Hình vuông có diện tích $16\text{ cm}^2$ thì chu vi của nó là $8\text{ cm}$.}
{\True Trong các hình chữ nhật có cùng chu vi thì hình vuông có diện tích lớn nhất.}
{Diện tích hình chữ nhật bằng (dài + rộng) nhân 2.}
\loigiai{
\begin{itemize}
    \item Ý a) Đúng: Quy chuẩn đo lường.
    \item Ý b) Sai: Sai vì cạnh là $4\text{ cm}$, chu vi là $4 \times 4 = 16\text{ cm}$.
    \item Ý c) Đúng: Định lí cực trị hình học.
    \item Ý d) Sai: Sai vì đó là công thức tính chu vi.
\end{itemize}
\begin{itemize}
    \item a) \textbf{Đúng}.
    \item b) \textbf{Sai}.
    \item c) \textbf{Đúng}.
    \item d) \textbf{Sai}.
\end{itemize}
}
\end{ex}

\begin{ex} %[3H3V53]
Cho hình vuông $ABCD$ có cạnh dài $8\text{ cm}$:
\choiceTF
{\True Chu vi hình vuông $ABCD$ bằng $32\text{ cm}$.}
{\True Diện tích hình vuông $ABCD$ bằng $64\text{ cm}^2$.}
{Nếu kéo dài một cạnh thêm $2\text{ cm}$ thì được hình chữ nhật có chu vi $34\text{ cm}$.}
{\True Có thể chia hình vuông đó thành 4 hình vuông nhỏ bằng nhau có cạnh $4\text{ cm}$.}
\loigiai{
\begin{itemize}
    \item Ý a) Đúng: $8 \times 4 = 32\text{ cm}$.
    \item Ý b) Đúng: $8 \times 8 = 64\text{ cm}^2$.
    \item Ý c) Sai: Sai vì chu vi tăng thêm $2 \times 2 = 4\text{ cm}$ thành $36\text{ cm}$.
    \item Ý d) Đúng: $8 : 4 = 2$ phần mỗi chiều.
\end{itemize}
\begin{itemize}
    \item a) \textbf{Đúng}.
    \item b) \textbf{Đúng}.
    \item c) \textbf{Sai}.
    \item d) \textbf{Đúng}.
\end{itemize}
}
\end{ex}

\textbf{Phần 3. Trả lời ngắn}
\setcounterpageref{ex}{0}

\begin{ex} %[3H3V53]
Một hình vuông có chu vi $44\text{ cm}$. Tính cạnh hình vuông theo xăng-ti-mét.
\shortans{11}
\loigiai{
\textbf{Đáp số:} $11$\\
$44 : 4 = 11\text{ (cm)}$.
}
\end{ex}

\begin{ex} %[3H3V53]
Một hình chữ nhật có chiều dài $18\text{ cm}$, chiều rộng $5\text{ cm}$. Tính diện tích hình chữ nhật theo xăng-ti-mét vuông.
\shortans{90}
\loigiai{
\textbf{Đáp số:} $90$\\
$18 \times 5 = 90\text{ (cm}^2)$.
}
\end{ex}

\begin{ex} %[3H3V53]
Hình vuông có diện tích $49\text{ cm}^2$. Tính chu vi hình vuông đó theo xăng-ti-mét.
\shortans{28}
\loigiai{
\textbf{Đáp số:} $28$\\
Cạnh là $7\text{ cm}$. Chu vi: $7 \times 4 = 28\text{ (cm)}$.
}
\end{ex}

\begin{ex} %[3H3V53]
Một hình chữ nhật có chu vi $36\text{ cm}$, chiều dài hơn chiều rộng $4\text{ cm}$. Tính diện tích hình chữ nhật theo xăng-ti-mét vuông.
\shortans{77}
\loigiai{
\textbf{Đáp số:} $77$\\
Nửa chu vi: $18\text{ cm}$. Dài: $11\text{ cm}$, rộng: $7\text{ cm}$. Diện tích: $11 \times 7 = 77\text{ (cm}^2)$.
}
\end{ex}

\begin{ex} %[3H3V53]
Cần bao nhiêu mét dây thép gai để rào kín một mảnh vườn hình vuông cạnh $22\text{ m}$ nếu chừa cổng ra vào rộng $2\text{ m}$?
\shortans{86}
\loigiai{
\textbf{Đáp số:} $86$\\
Chu vi: $22 \times 4 = 88\text{ m}$. Cần: $88 - 2 = 86\text{ (m)}$.
}
\end{ex}

\begin{ex} %[3H3V53]
Một hình chữ nhật có diện tích $60\text{ cm}^2$, chiều rộng $6\text{ cm}$. Tính chu vi hình chữ nhật đó theo xăng-ti-mét.
\shortans{32}
\loigiai{
\textbf{Đáp số:} $32$\\
Chiều dài: $60 : 6 = 10\text{ cm}$. Chu vi: $(10 + 6) \times 2 = 32\text{ (cm)}$.
}
\end{ex}
